\section{Cohomology of affine schemes}
\label{section:III.1}

\subsection{Review of the exterior algebra complex}
\label{subsection:III.1.1}

\begin{env}[1.1.1]
\label{III.1.1.1}
Let $A$ be a ring and $\mathbf f=(f_i)_{1\leq i\leq r}$ a system of $r$
elements of $A$. The \emph{exterior algebra complex
$K_\bullet(\mathbf f)$} corresponding to $\mathbf f$ is a complex of chains
(G, I, 2.2) defined in the following way: the graded $A$-module
$K_\bullet(\mathbf f)$ is equal to the \emph{exterior algebra
$\bigwedge(A^r)$}, graded in the usual way, and the boundary operator is the
\emph{interior multiplication $i_{\mathbf f}$} by $\mathbf f$ regarded as
an element of the dual $(A^r)^\vee$; recall that $i_{\mathbf f}$ is an
\emph{antiderivation} of degree $-1$ of $\bigwedge(A^r)$, and that if
$(\mathbf e_i)_{1\leq i\leq r}$ is the canonical basis of $A^r$, we have
$i_{\mathbf f}(\mathbf e_i)=f_i$; the verification of the condition
$i_{\mathbf f}\circ i_{\mathbf f}=0$ is immediate.

An equivalent definition is the following: for each $i$, we consider
a complex of chains $K_\bullet(f_i)$ defined as follows:
$K_0(f_i)=K_1(f_i)=A$, $K_n(f_i)=0$ for $n\neq0,1$; the boundary operator is
defined by the condition that $d_1:A\to A$ is \emph{multiplication by
$f_i$}. We then take for $K_\bullet(\mathbf f)$ the \emph{tensor product
$K_\bullet(f_1)\otimes K_\bullet(f_2)\otimes\cdots\otimes K_\bullet(f_r)$}
(G, I, 2.7) equipped with its total degree; the verification of the
isomorphism of this complex with the complex defined above is immediate.
\end{env}

\begin{env}[1.1.2]
\label{III.1.1.2}
For every $A$-module $M$, we define the \emph{complex of chains}
\begin{equation}
\label{III.1.1.2.1}
\tag{1.1.2.1}
K_\bullet(\mathbf f,M)=K_\bullet(\mathbf f)\otimes_A M
\end{equation}
and the \emph{complex of cochains} (G, I, 2.2)
\begin{equation}
\label{III.1.1.2.2}
\tag{1.1.2.2}
K^\bullet(\mathbf f,M)=\operatorname{Hom}_A(K_\bullet(\mathbf f),M).
\end{equation}
If $g$ is a $k$-cochain of the latter complex, and if we set
\[
g(i_1,\ldots,i_k)=g(\mathbf e_{i_1}\wedge\cdots\wedge\mathbf e_{i_k}),
\]
$g$ is identified with an \emph{alternating} map from $[1,r]^k$ to $M$, and
it follows from the preceding definitions that we have
\begin{equation}
\label{III.1.1.2.3}
\tag{1.1.2.3}
d^kg(i_1,i_2,\ldots,i_{k+1})=
\sum_{h=1}^{k+1}(-1)^{h-1}f_{i_h}
g(i_1,\ldots,\widehat{i_h},\ldots,i_{k+1}).
\end{equation}
\end{env}

\oldpage[III]{83}
\begin{env}[1.1.3]
\label{III.1.1.3}
From the preceding complexes we deduce, as usual, the \emph{homology and
cohomology $A$-modules} (G, I, 2.2)
\begin{align}
\label{III.1.1.3.1}
\tag{1.1.3.1}
H_\bullet(\mathbf f,M)&=H_\bullet(K_\bullet(\mathbf f,M)),\\
\label{III.1.1.3.2}
\tag{1.1.3.2}
H^\bullet(\mathbf f,M)&=H^\bullet(K^\bullet(\mathbf f,M)).
\end{align}
Moreover, we define an $A$-isomorphism
$K_i(\mathbf f,M)\xrightarrow{\sim}K^{r-i}(\mathbf f,M)$ by assigning to
every chain
$z=\sum(\mathbf e_{i_1}\wedge\cdots\wedge\mathbf e_{i_k})\otimes
z_{i_1\ldots i_k}$ the cochain $g_z$ such that
$g_z(j_1,\ldots,j_{r-k})=\varepsilon z_{i_1\ldots i_k}$, where
$(j_h)_{1\leq h\leq r-k}$ is the strictly increasing sequence complementary
to the strictly increasing sequence $(i_h)_{1\leq h\leq k}$ in $[1,r]$ and
$\varepsilon=(-1)^\nu$, $\nu$ being the number of inversions of the
permutation $i_1,\ldots,i_k,j_1,\ldots,j_{r-k}$ of $[1,r]$. One verifies that
$g_{dz}=d(g_z)$, which gives an isomorphism
\begin{equation}
\label{III.1.1.3.3}
\tag{1.1.3.3}
H^i(\mathbf f,M)\xrightarrow{\sim}H_{r-i}(\mathbf f,M)
\qquad\text{for }0\leq i\leq r.
\end{equation}
In this chapter, we will mainly have to consider the cohomology modules
$H^\bullet(\mathbf f,M)$.

For a given $\mathbf f$, it is immediate (G, I, 2.1) that
$M\mapsto H^\bullet(\mathbf f,M)$ is a \emph{cohomological functor}
(T, II, 2.1) from the category of $A$-modules to that of graded
$A$-modules, zero in degrees $<0$ and $>r$. Moreover, we have
\begin{equation}
\label{III.1.1.3.4}
\tag{1.1.3.4}
H^0(\mathbf f,M)=\operatorname{Hom}_A(A/(\mathbf f),M)
\end{equation}
where $(\mathbf f)$ denotes the ideal of $A$ generated by
$f_1,\ldots,f_r$; this follows immediately from
(\hyperref[III.1.1.2.3]{1.1.2.3}), and it is clear that
$H^0(\mathbf f,M)$ is identified with the submodule of $M$ \emph{annihilated
by $(\mathbf f)$}. Similarly, by
(\hyperref[III.1.1.2.3]{1.1.2.3}) we have
\begin{equation}
\label{III.1.1.3.5}
\tag{1.1.3.5}
H^r(\mathbf f,M)=M/\Bigl(\sum_{i=1}^r f_iM\Bigr)
=(A/(\mathbf f))\otimes_A M.
\end{equation}
We will use the following known result, of which we will recall a
proof for the sake of completeness:
\end{env}

\begin{proposition}[1.1.4]
\label{III.1.1.4}
Let $A$ be a ring, $\mathbf f=(f_i)_{1\leq i\leq r}$ a finite family
of elements of $A$, and $M$ an $A$-module. If, for $1\leq i\leq r$, the
homothety $z\mapsto f_i z$ in $M_{i-1}=M/(f_1M+\cdots+f_{i-1}M)$ is injective,
then $H^i(\mathbf f,M)=0$ for $i\neq r$.
\end{proposition}

\begin{proof}
Indeed, by virtue of (\hyperref[III.1.1.3.3]{1.1.3.3}), it suffices to
prove that $H_i(\mathbf f,M)=0$ for every $i>0$. We argue by induction on
$r$, the case $r=0$ being trivial. Set
$\mathbf f'=(f_i)_{1\leq i\leq r-1}$; this family satisfies the conditions of
the statement, so, if we set $L_\bullet=K_\bullet(\mathbf f',M)$, we have
$H_i(L_\bullet)=0$ for $i>0$ by hypothesis, and
$H_0(L_\bullet)=M_{r-1}$ by virtue of
(\hyperref[III.1.1.3.3]{1.1.3.3}) and
(\hyperref[III.1.1.3.5]{1.1.3.5}). For brevity, set
$K_\bullet=K_\bullet(f_r)=K_0\oplus K_1$, with $K_0=K_1=A$,
$d_1:K_1\to K_0$ being multiplication by $f_r$; by definition
(\hyperref[III.1.1.1]{1.1.1}) we have
$K_\bullet(\mathbf f,M)=K_\bullet\otimes_A L_\bullet$. Now, we have the
following lemma:

\begin{lemma}[1.1.4.1]
\label{III.1.1.4.1}
Let $K_\bullet$ be a complex of chains formed of free $A$-modules that are
zero except in dimensions $0$ and $1$. For every complex of chains $L_\bullet$
formed of $A$-modules, we have an exact sequence
\[
0\longrightarrow H_0(K_\bullet\otimes H_p(L_\bullet))
\longrightarrow H_p(K_\bullet\otimes L_\bullet)
\longrightarrow H_1(K_\bullet\otimes H_{p-1}(L_\bullet))
\longrightarrow0
\]
for every index $p$.
\end{lemma}

\oldpage[III]{84}
This is a particular case of an exact sequence of low-degree terms of the
Künneth spectral sequence (M, XVII, 5.2 a) and G, I, 5.5.2); it can be
proved directly in the following way. Let us regard $K_0$ and $K_1$ as
complexes of chains (zero in dimensions $\neq0$ and $\neq1$
respectively); we then have an exact sequence of complexes
\[
0\longrightarrow K_0\otimes L_\bullet
\longrightarrow K_\bullet\otimes L_\bullet
\longrightarrow K_1\otimes L_\bullet\longrightarrow0
\]
to which we can apply the homology exact sequence
\[
\cdots\longrightarrow H_{p+1}(K_1\otimes L_\bullet)
\xrightarrow{\partial}H_p(K_0\otimes L_\bullet)
\longrightarrow H_p(K_\bullet\otimes L_\bullet)
\longrightarrow H_p(K_1\otimes L_\bullet)
\xrightarrow{\partial}H_{p-1}(K_0\otimes L_\bullet)
\longrightarrow\cdots.
\]
But it is evident that
$H_p(K_0\otimes L_\bullet)=K_0\otimes H_p(L_\bullet)$ and
$H_p(K_1\otimes L_\bullet)=K_1\otimes H_{p-1}(L_\bullet)$ for every $p$;
moreover, one verifies immediately that the operator
$\partial:K_1\otimes H_p(L_\bullet)\to K_0\otimes H_p(L_\bullet)$ is none
other than $d_1\otimes1$; the lemma therefore follows from the preceding exact
sequence and from the definition of $H_0(K_\bullet\otimes H_p(L_\bullet))$ and
of $H_1(K_\bullet\otimes H_{p-1}(L_\bullet))$.

This lemma being established, the end of the proof of
(\hyperref[III.1.1.4]{1.1.4}) is immediate: the induction hypothesis and
this lemma (\hyperref[III.1.1.4.1]{1.1.4.1}) give
$H_p(K_\bullet\otimes L_\bullet)=0$ for $p\geq2$; moreover, if we prove that
$H_1(K_\bullet,H_0(L_\bullet))=0$, we will also deduce
$H_1(K_\bullet\otimes L_\bullet)=0$ from Lemma
(\hyperref[III.1.1.4.1]{1.1.4.1}); but by definition,
$H_1(K_\bullet,H_0(L_\bullet))$ is none other than the kernel of the
homothety $z\mapsto f_rz$ in $M_{r-1}$, and since this kernel is zero by
hypothesis, this completes the proof.
\end{proof}

\begin{env}[1.1.5]
\label{III.1.1.5}
Let $\mathbf g=(g_i)_{1\leq i\leq r}$ be a second sequence of $r$ elements of
$A$, and set $\mathbf f\mathbf g=(f_ig_i)_{1\leq i\leq r}$. We can define
a canonical homomorphism of complexes
\begin{equation}
\label{III.1.1.5.1}
\tag{1.1.5.1}
\varphi_{\mathbf g}:K_\bullet(\mathbf f\mathbf g)\longrightarrow
K_\bullet(\mathbf f)
\end{equation}
as the canonical extension to the exterior algebra $\bigwedge(A^r)$ of
the $A$-linear map
$(x_1,\ldots,x_r)\mapsto(g_1x_1,\ldots,g_rx_r)$ of $A^r$ into itself. To
see that this is indeed a homomorphism of complexes, it suffices to remark,
in general, that if $u:E\to F$ is an $A$-linear map,
$\mathbf x\in F^\vee$ and $\mathbf y={}^tu(\mathbf x)\in E^\vee$, then we have
the formula
\begin{equation}
\label{III.1.1.5.2}
\tag{1.1.5.2}
(\bigwedge u)\circ i_{\mathbf y}=i_{\mathbf x}\circ(\bigwedge u)~;
\end{equation}
indeed, the two sides are antiderivations of $\bigwedge F$, and it
suffices to verify that they coincide in $F$, which follows immediately from
the definitions.

When $K_\bullet(\mathbf f)$ is identified with the tensor product of the
$K_\bullet(f_i)$ (\hyperref[III.1.1.1]{1.1.1}), $\varphi_{\mathbf g}$ is
the tensor product of the $\varphi_{g_i}$, where $\varphi_{g_i}$ reduces to
the identity in degree $0$ and to multiplication by $g_i$ in degree $1$.
\end{env}

\begin{env}[1.1.6]
\label{III.1.1.6}
In particular, for every pair of integers $m,n$ such that $0\leq n\leq m$,
we have homomorphisms of complexes
\begin{equation}
\label{III.1.1.6.1}
\tag{1.1.6.1}
\varphi_{\mathbf f^{m-n}}:K_\bullet(\mathbf f^m)\longrightarrow
K_\bullet(\mathbf f^n)
\end{equation}
and consequently homomorphisms
\begin{align}
\label{III.1.1.6.2}
\tag{1.1.6.2}
\varphi_{\mathbf f^{m-n}} &:K^\bullet(\mathbf f^n,M)\longrightarrow
K^\bullet(\mathbf f^m,M),\\
\label{III.1.1.6.3}
\tag{1.1.6.3}
\varphi_{\mathbf f^{m-n}} &:H^\bullet(\mathbf f^n,M)\longrightarrow
H^\bullet(\mathbf f^m,M).
\end{align}

\oldpage[III]{85}
The latter evidently satisfy the transitivity condition
$\varphi_{\mathbf f^{m-p}}=
\varphi_{\mathbf f^{m-n}}\circ\varphi_{\mathbf f^{n-p}}$ for $p\leq n\leq m$;
they therefore define two \emph{inductive systems} of $A$-modules; we set
\begin{align}
\label{III.1.1.6.4}
\tag{1.1.6.4}
C^\bullet((\mathbf f),M)&=\varinjlim_n K^\bullet(\mathbf f^n,M),\\
\label{III.1.1.6.5}
\tag{1.1.6.5}
H^\bullet((\mathbf f),M)&=H^\bullet(C^\bullet((\mathbf f),M))
=\varinjlim_n H^\bullet(\mathbf f^n,M),
\end{align}
the last equality coming from the fact that passage to the inductive limit
commutes with the functor $H^\bullet$ (G, I, 2.1). We will see later
(\hyperref[III.1.4.3]{1.4.3}) that $H^\bullet((\mathbf f),M)$ in fact
depends only on the \emph{ideal $(\mathbf f)$} in $A$ (and even only on the
$(\mathbf f)$-preadic topology on $A$), which justifies the notation.

It is clear that $M\mapsto C^\bullet((\mathbf f),M)$ is an exact
$A$-linear functor, and $M\mapsto H^\bullet((\mathbf f),M)$ a
cohomological functor.
\end{env}

\begin{env}[1.1.7]
\label{III.1.1.7}
Let $\mathbf f=(f_i)\in A^r$, $\mathbf g=(g_i)\in A^r$; denote by
$e_{\mathbf g}$ left multiplication by the vector
$\mathbf g\in A^r$ in the exterior algebra $\bigwedge(A^r)$; we know that
we have the \emph{homotopy formula}
\begin{equation}
\label{III.1.1.7.1}
\tag{1.1.7.1}
i_{\mathbf f}e_{\mathbf g}+e_{\mathbf g}i_{\mathbf f}
=\langle\mathbf g,\mathbf f\rangle 1
\end{equation}
in the $A$-module $\bigwedge(A^r)$ ($1$ denoting the identity automorphism
of $\bigwedge(A^r)$);
this relation also means that, \emph{in the complex
$K_\bullet(\mathbf f)$}, we have
\begin{equation}
\label{III.1.1.7.2}
\tag{1.1.7.2}
de_{\mathbf g}+e_{\mathbf g}d=\langle\mathbf g,\mathbf f\rangle 1.
\end{equation}
If the ideal $(\mathbf f)$ is equal to $A$, there exists $\mathbf g\in A^r$
such that $\langle\mathbf g,\mathbf f\rangle=\sum_{i=1}^r g_if_i=1$.
Consequently (G, I, 2.4):
\end{env}

\begin{proposition}[1.1.8]
\label{III.1.1.8}
Suppose that the ideal $(\mathbf f)$ generated by the $f_i$ is equal to $A$.
Then the complex $K_\bullet(\mathbf f)$ is homotopically trivial, and
therefore so are the complexes $K_\bullet(\mathbf f,M)$ and
$K^\bullet(\mathbf f,M)$ for every $A$-module $M$.
\end{proposition}

\begin{corollary}[1.1.9]
\label{III.1.1.9}
If $(\mathbf f)=A$, we have $H^\bullet(\mathbf f,M)=0$ and
$H^\bullet((\mathbf f),M)=0$ for every $A$-module $M$.
\end{corollary}

\begin{proof}
Indeed, we then have $(\mathbf f^n)=A$ for every $n$.
\end{proof}

\begin{remark}[1.1.10]
\label{III.1.1.10}
With the same notation as above, let $X=\operatorname{Spec}(A)$, and
let $Y$ be the closed subscheme of $X$ defined by the ideal $(\mathbf f)$. We
will prove in \S~9 that $H^\bullet((\mathbf f),M)$ is isomorphic to the
cohomology $H_Y^\bullet(X,\widetilde M)$ corresponding to the antifilter
$\Phi$ of closed subsets of $Y$ (T, 3.2). We will also show that
Proposition (\hyperref[III.1.2.3]{1.2.3}), applied to $X$ and to
$\mathcal F=\widetilde M$, is a particular case of a cohomology exact
sequence
\[
\cdots\longrightarrow H_Y^p(X,\mathcal F)\longrightarrow H^p(X,\mathcal F)
\longrightarrow H^p(X-Y,\mathcal F)\longrightarrow H_Y^{p+1}(X,\mathcal F)
\longrightarrow\cdots.
\]
\end{remark}

\subsection{Čech cohomology of an open covering}
\label{subsection:III.1.2}

\begin{notation}[1.2.1]
\label{III.1.2.1}
In this subsection, we denote by:
\begin{itemize}
\item $X$ a prescheme;
\item $\mathcal F$ a quasi-coherent $\mathcal O_X$-module;
\oldpage[III]{86}
\item $A=\Gamma(X,\mathcal O_X)$, $M=\Gamma(X,\mathcal F)$;
\item $\mathbf f=(f_i)_{1\leq i\leq r}$ a finite system of elements of $A$;
\item $U_i=X_{f_i}$ the open set
(\hyperref[0.5.5.2]{0\textsubscript{I}, 5.5.2}) of the $x\in X$ such that
$f_i(x)\neq0$;
\item $U=\bigcup_{i=1}^r U_i$;
\item $\mathfrak U$ the covering $(U_i)_{1\leq i\leq r}$ of $U$.
\end{itemize}
\end{notation}

\begin{env}[1.2.2]
\label{III.1.2.2}
Suppose that $X$ is either a prescheme whose underlying space is
\emph{Noetherian}, or a \emph{scheme} whose underlying space is
\emph{quasi-compact}. We then know (\hyperref[I.9.3.3]{I, 9.3.3}) that
$\Gamma(U_i,\mathcal F)=M_{f_i}$. We set
\[
U_{i_0i_1\ldots i_p}=\bigcap_{k=0}^p U_{i_k}
=X_{f_{i_0}f_{i_1}\cdots f_{i_p}}
\]
(\hyperref[0.5.5.3]{0\textsubscript{I}, 5.5.3}); hence we also have
\begin{equation}
\label{III.1.2.2.1}
\tag{1.2.2.1}
\Gamma(U_{i_0i_1\ldots i_p},\mathcal F)=M_{f_{i_0}f_{i_1}\cdots f_{i_p}}.
\end{equation}
Now (\hyperref[0.1.6.1]{0\textsubscript{I}, 1.6.1})
$M_{f_{i_0}f_{i_1}\cdots f_{i_p}}$ is identified with the inductive limit
$\varinjlim_n M_{i_0i_1\ldots i_p}^{(n)}$, where the inductive system is
formed by the $M_{i_0i_1\ldots i_p}^{(n)}=M$, the homomorphism
$\varphi_{nm}:M_{i_0i_1\ldots i_p}^{(m)}\to
M_{i_0i_1\ldots i_p}^{(n)}$ being multiplication by
$(f_{i_0}f_{i_1}\cdots f_{i_p})^{n-m}$ for $m\leq n$. Denote by
$C_n^p(M)$ the set of \emph{alternating} maps from $[1,r]^{p+1}$
to $M$ (for every $n$); these $A$-modules again form an inductive system
for the $\varphi_{nm}$. If $C^p(\mathfrak U,\mathcal F)$ is the group of
\emph{alternating} Čech $p$-cochains relative to the covering $\mathfrak U$,
with coefficients in $\mathcal F$ (G, II, 5.1), it follows from the above
that we can write
\begin{equation}
\label{III.1.2.2.2}
\tag{1.2.2.2}
C^p(\mathfrak U,\mathcal F)=\varinjlim_n C_n^p(M).
\end{equation}
Now, with the notation of (\hyperref[III.1.1.2]{1.1.2}), $C_n^p(M)$
is identified with $K^{p+1}(\mathbf f^n,M)$, and the map $\varphi_{nm}$
is identified with the map $\varphi_{\mathbf f^{n-m}}$ defined in
(\hyperref[III.1.1.6]{1.1.6}). We therefore have, for every $p\geq0$, a
canonical isomorphism, functorial in $\mathcal F$,
\begin{equation}
\label{III.1.2.2.3}
\tag{1.2.2.3}
C^p(\mathfrak U,\mathcal F)\xrightarrow{\sim}C^{p+1}((\mathbf f),M).
\end{equation}
Moreover, formula (\hyperref[III.1.1.2.3]{1.1.2.3}) and the definition of
the cohomology of a covering (G, II, 5.1) show that the isomorphisms
(\hyperref[III.1.2.2.3]{1.2.2.3}) are compatible with the coboundary
operators.
\end{env}

\begin{proposition}[1.2.3]
\label{III.1.2.3}
If $X$ is a prescheme whose underlying space is Noetherian, or a
scheme whose underlying space is quasi-compact, there exists a canonical
isomorphism, functorial in $\mathcal F$,
\begin{equation}
\label{III.1.2.3.1}
\tag{1.2.3.1}
H^p(\mathfrak U,\mathcal F)\xrightarrow{\sim}H^{p+1}((\mathbf f),M)
\qquad\text{for every }p\geq1.
\end{equation}
Moreover, we have an exact sequence, functorial in $\mathcal F$,
\begin{equation}
\label{III.1.2.3.2}
\tag{1.2.3.2}
0\longrightarrow H^0((\mathbf f),M)\longrightarrow M
\longrightarrow H^0(\mathfrak U,\mathcal F)
\longrightarrow H^1((\mathbf f),M)\longrightarrow0.
\end{equation}
\end{proposition}

\oldpage[III]{87}
\begin{proof}
Indeed, the relations (\hyperref[III.1.2.3.1]{1.2.3.1}) are immediate
consequences of what we saw in (\hyperref[III.1.2.2]{1.2.2}). On the
other hand, we have
$C^0(\mathfrak U,\mathcal F)=C^1((\mathbf f),M)$;
$H^0(\mathfrak U,\mathcal F)$ is consequently identified with the subgroup of
$1$-cocycles of $C^1((\mathbf f),M)$; since
$M=C^0((\mathbf f),M)$, the exact sequence
(\hyperref[III.1.2.3.2]{1.2.3.2}) is none other than the one that results
from the definition of the cohomology groups $H^0((\mathbf f),M)$ and
$H^1((\mathbf f),M)$.
\end{proof}

\begin{corollary}[1.2.4]
\label{III.1.2.4}
Suppose that the $X_{f_i}$ are quasi-compact and that there exist
$g_i\in\Gamma(U,\mathcal O_X)$ such that
$\sum_i g_i(f_i\vert U)=1\vert U$. Then, for every quasi-coherent
$(\mathcal O_X\vert U)$-module $\mathcal G$, we have
$H^p(\mathfrak U,\mathcal G)=0$ for $p>0$; if moreover $U=X$,
the canonical homomorphism (\hyperref[III.1.2.3.2]{1.2.3.2})
$\Gamma(X,\mathcal G)\to H^0(\mathfrak U,\mathcal G)$ is bijective.
\end{corollary}

\begin{proof}
Since the $U_i=X_{f_i}$ are quasi-compact by hypothesis, so is $U$, and we
may therefore restrict to the case where $U=X$; the hypothesis then implies
$H^p((\mathbf f),M)=0$ for every $p\geq0$
(\hyperref[III.1.1.9]{1.1.9}). The corollary then follows immediately from
(\hyperref[III.1.2.3.1]{1.2.3.1}) and
(\hyperref[III.1.2.3.2]{1.2.3.2}).
\end{proof}

Observe that, since
$H^0(\mathfrak U,\mathcal F)=H^0(U,\mathcal F)$ (G, II, 5.2.2), we have thus
proved (I, 1.3.7) again as a particular case.

\begin{remark}[1.2.5]
\label{III.1.2.5}
Suppose that $X$ is an \emph{affine scheme}; then the
$U_i=X_{f_i}=D(f_i)$ are affine open subsets, as are the
$U_{i_0i_1\ldots i_p}$ (but $U$ is not necessarily affine). In this
case, the functors $\Gamma(X,\mathcal F)$ and
$\Gamma(U_{i_0i_1\ldots i_p},\mathcal F)$ are exact in $\mathcal F$
(I, 1.3.11). If we have an exact sequence
$0\to\mathcal F'\to\mathcal F\to\mathcal F''\to0$ of quasi-coherent
$\mathcal O_X$-modules, the sequence of complexes
\[
0\longrightarrow C^\bullet(\mathfrak U,\mathcal F')
\longrightarrow C^\bullet(\mathfrak U,\mathcal F)
\longrightarrow C^\bullet(\mathfrak U,\mathcal F'')
\longrightarrow0
\]
is exact, and therefore gives a cohomology exact sequence
\[
\cdots\longrightarrow H^p(\mathfrak U,\mathcal F')
\longrightarrow H^p(\mathfrak U,\mathcal F)
\longrightarrow H^p(\mathfrak U,\mathcal F'')
\xrightarrow{\partial}H^{p+1}(\mathfrak U,\mathcal F')
\longrightarrow\cdots.
\]
On the other hand, if we set $M'=\Gamma(X,\mathcal F')$,
$M''=\Gamma(X,\mathcal F'')$, the sequence
$0\to M'\to M\to M''\to0$ is exact; since
$C^\bullet((\mathbf f),M)$ is an exact functor in $M$, we also have the
cohomology exact sequence
\[
\cdots\longrightarrow H^p((\mathbf f),M')
\longrightarrow H^p((\mathbf f),M)
\longrightarrow H^p((\mathbf f),M'')
\xrightarrow{\partial}H^{p+1}((\mathbf f),M')
\longrightarrow\cdots.
\]
This being so, since the diagram
\[
\xymatrix{
0\ar[r]&C^\bullet(\mathfrak U,\mathcal F')\ar[r]\ar[d]
&C^\bullet(\mathfrak U,\mathcal F)\ar[r]\ar[d]
&C^\bullet(\mathfrak U,\mathcal F'')\ar[r]\ar[d]&0\\
0\ar[r]&C^\bullet((\mathbf f),M')\ar[r]
&C^\bullet((\mathbf f),M)\ar[r]
&C^\bullet((\mathbf f),M'')\ar[r]&0
}
\]
\oldpage[III]{88}
is commutative, we conclude that the diagrams
\begin{equation}
\label{III.1.2.5.1}
\tag{1.2.5.1}
\xymatrix{
H^p(\mathfrak U,\mathcal F'')\ar[r]^{\partial}\ar[d]
&H^{p+1}(\mathfrak U,\mathcal F')\ar[d]\\
H^{p+1}((\mathbf f),M'')\ar[r]^{\partial}
&H^{p+2}((\mathbf f),M')
}
\end{equation}
are commutative for every $p$ (G, I, 2.1.1).
\end{remark}

\subsection{Cohomology of an affine scheme}
\label{subsection:III.1.3}

\begin{theorem}[1.3.1]
\label{III.1.3.1}
Let $X$ be an affine scheme. For every quasi-coherent $\mathcal O_X$-module
$\mathcal F$, we have $H^p(X,\mathcal F)=0$ for every $p>0$.
\end{theorem}

\begin{proof}
Let $\mathfrak U$ be a finite covering of $X$ by affine open subsets
$X_{f_i}=D(f_i)$ ($1\leq i\leq r$); we know that then the ideal generated in
$A=\Gamma(X,\mathcal O_X)$ by the $f_i$ is equal to $A$. We therefore
conclude from (\hyperref[III.1.2.4]{1.2.4}) that
$H^p(\mathfrak U,\mathcal F)=0$ for $p>0$. Since there are finite
coverings of $X$ by affine open subsets that are arbitrarily fine
(I, 1.1.10), the definition of Čech cohomology (G, II, 5.8) shows that we
also have $\check H^p(X,\mathcal F)=0$ for $p>0$. But this also applies to
every prescheme $X_f$ for $f\in A$ (I, 1.3.6), so
$\check H^p(X_f,\mathcal F)=0$ for $p>0$. Since
$X_f\cap X_g=X_{fg}$, we deduce that we also have
$H^p(X,\mathcal F)=0$ for every $p>0$, by virtue of (G, II, 5.9.2).
\end{proof}

\begin{corollary}[1.3.2]
\label{III.1.3.2}
Let $Y$ be a prescheme and $f:X\to Y$ an affine morphism (II, 1.6.1). For
every quasi-coherent $\mathcal O_X$-module $\mathcal F$, we have
$R^qf_*(\mathcal F)=0$ for $q>0$.
\end{corollary}

\begin{proof}
Indeed, by definition $R^qf_*(\mathcal F)$ is the $\mathcal O_Y$-module
associated to the presheaf $U\mapsto H^q(f^{-1}(U),\mathcal F)$, $U$
ranging over the open subsets of $Y$. But the affine open subsets form a
basis of $Y$, and for such an open subset $U$, $f^{-1}(U)$ is affine
(II, 1.3.2), so $H^q(f^{-1}(U),\mathcal F)=0$ by
(\hyperref[III.1.3.1]{1.3.1}), which proves the corollary.
\end{proof}

\begin{corollary}[1.3.3]
\label{III.1.3.3}
Let $Y$ be a prescheme and $f:X\to Y$ an affine morphism. For every
quasi-coherent $\mathcal O_X$-module $\mathcal F$, the canonical homomorphism
$H^p(Y,f_*(\mathcal F))\to H^p(X,\mathcal F)$ (0, 12.1.3.1) is bijective
for every $p$.
\end{corollary}

\begin{proof}
Indeed, it suffices (by virtue of (0, 12.1.7)) to show that the
edge homomorphisms
$''E_2^{p0}=H^p(Y,f_*(\mathcal F))\to H^p(X,\mathcal F)$ of the second
spectral sequence of the composite functor $\Gamma f_*$ are bijective. But
the $E_2$ term of this sequence is given by
$''E_2^{pq}=H^p(Y,R^qf_*(\mathcal F))$ (G, II, 4.17.1), so it follows from
(\hyperref[III.1.3.2]{1.3.2}) that $''E_2^{pq}=0$ for $q>0$, and the
spectral sequence degenerates; whence our assertion (0, 11.1.6).
\end{proof}

\begin{corollary}[1.3.4]
\label{III.1.3.4}
Let $f:X\to Y$ be an affine morphism and $g:Y\to Z$ a morphism.
\oldpage[III]{89}
For every quasi-coherent $\mathcal O_X$-module $\mathcal F$,
the canonical homomorphism
$R^pg_*(f_*(\mathcal F))\to R^p(g\circ f)_*(\mathcal F)$ (0, 12.2.5.1)
is bijective for every $p$.
\end{corollary}

\begin{proof}
It suffices to remark that, by
(\hyperref[III.1.3.3]{1.3.3}), for every affine open subset $W$ of $Z$,
the canonical homomorphism
$H^p(g^{-1}(W),f_*(\mathcal F))\to
H^p(f^{-1}(g^{-1}(W)),\mathcal F)$ is bijective; this proves that
the homomorphism of presheaves defining the canonical homomorphism
$R^pg_*(f_*(\mathcal F))\to R^p(g\circ f)_*(\mathcal F)$ is bijective
(0, 12.2.5).
\end{proof}


\subsection{Application to the cohomology of arbitrary preschemes}
\label{subsection:III.1.4}

\begin{proposition}[1.4.1]
\label{III.1.4.1}
Let $X$ be a scheme and $\mathfrak U=(U_\alpha)$ a covering of $X$ by
affine open subsets. For every quasi-coherent $\mathcal O_X$-module
$\mathcal F$, the cohomology modules $H^\bullet(X,\mathcal F)$ and
$H^\bullet(\mathfrak U,\mathcal F)$ (over $\Gamma(X,\mathcal O_X)$) are
canonically isomorphic.
\end{proposition}

\begin{proof}
Indeed, since $X$ is a scheme, every finite intersection $V$ of open subsets
of the covering $\mathfrak U$ is affine (I, 5.5.6), so
$H^q(V,\mathcal F)=0$ for $q\geq1$ by virtue of
(\hyperref[III.1.3.1]{1.3.1}). The proposition then follows from Leray's
theorem (G, II, 5.4.1).
\end{proof}

\begin{remark}[1.4.2]
\label{III.1.4.2}
Note that the conclusion of (\hyperref[III.1.4.1]{1.4.1}) is still
valid when the finite intersections of the sets $U_\alpha$ are
affine, even when $X$ is not necessarily assumed to be a
scheme.
\end{remark}

\begin{corollary}[1.4.3]
\label{III.1.4.3}
Let $X$ be a scheme whose underlying space is quasi-compact,
$A=\Gamma(X,\mathcal O_X)$, and $\mathbf f=(f_i)_{1\leq i\leq r}$ a finite
sequence of elements of $A$ such that the $X_{f_i}$ (notation of
(\hyperref[III.1.2.1]{1.2.1})) are affine. Then (with the notation
of (\hyperref[III.1.2.1]{1.2.1})), for every quasi-coherent
$\mathcal O_X$-module $\mathcal F$, we have a canonical isomorphism,
functorial in $\mathcal F$,
\begin{equation}
\label{III.1.4.3.1}
\tag{1.4.3.1}
H^q(U,\mathcal F)\xrightarrow{\sim}H^{q+1}((\mathbf f),M)
\qquad\text{for }q\geq1
\end{equation}
and an exact sequence, functorial in $\mathcal F$,
\begin{equation}
\label{III.1.4.3.2}
\tag{1.4.3.2}
0\longrightarrow H^0((\mathbf f),M)\longrightarrow M
\longrightarrow H^0(U,\mathcal F)
\longrightarrow H^1((\mathbf f),M)\longrightarrow0.
\end{equation}
\end{corollary}

\begin{proof}
Indeed, this follows immediately from
(\hyperref[III.1.4.1]{1.4.1}) and
(\hyperref[III.1.2.3]{1.2.3}).
\end{proof}

\begin{env}[1.4.4]
\label{III.1.4.4}
If $X$ is an \emph{affine scheme}, it follows from
(\hyperref[III.1.2.5]{1.2.5}) and
(\hyperref[III.1.4.1]{1.4.1}) that for every $q\geq0$, the diagrams
\begin{equation}
\label{III.1.4.4.1}
\tag{1.4.4.1}
\xymatrix{
H^q(U,\mathcal F'')\ar[r]^{\partial}\ar[d]
&H^{q+1}(U,\mathcal F')\ar[d]\\
H^{q+1}((\mathbf f),M'')\ar[r]^{\partial}
&H^{q+2}((\mathbf f),M')
}
\end{equation}
corresponding to an exact sequence
$0\to\mathcal F'\to\mathcal F\to\mathcal F''\to0$ of quasi-coherent
$\mathcal O_X$-modules (with the notation of
(\hyperref[III.1.2.5]{1.2.5})) are commutative.
\end{env}

\oldpage[III]{90}
\begin{proposition}[1.4.5]
\label{III.1.4.5}
Let $X$ be a quasi-compact scheme and $\mathcal L$ an invertible
$\mathcal O_X$-module, and consider the graded ring
$A_\bullet=\Gamma_\bullet(\mathcal L)$ (0\textsubscript{I}, 5.4.6); then
\[
H^\bullet(\mathcal F,\mathcal L)
=\bigoplus_{n\in\mathbb Z}H^\bullet(X,\mathcal F\otimes\mathcal L^{\otimes n})
\]
is a graded $A_\bullet$-module, and for every $f\in A_n$, we have a
canonical isomorphism
\begin{equation}
\label{III.1.4.5.1}
\tag{1.4.5.1}
H^\bullet(X_f,\mathcal F)
\xrightarrow{\sim}(H^\bullet(\mathcal F,\mathcal L))_{(f)}
\end{equation}
of $(A_\bullet)_{(f)}$-modules.
\end{proposition}

\begin{proof}
Since $X$ is a quasi-compact scheme, we can compute the cohomology of
all the $\mathcal O_X$-modules
$\mathcal F\otimes\mathcal L^{\otimes n}$ by means of one and the same
finite covering $\mathfrak U=(U_i)$ by affine open subsets such that the
restriction $\mathcal L\vert U_i$ is isomorphic to $\mathcal O_X\vert U_i$
for each $i$ (\hyperref[III.1.4.1]{1.4.1}). It is then immediate that the
$U_i\cap X_f$ are affine open subsets (I, 1.3.6), and we can therefore also
compute the cohomology
$H^\bullet(X_f,\mathcal F\otimes\mathcal L^{\otimes n})$ by means of the
covering $\mathfrak U\vert X_f=(U_i\cap X_f)$
(\hyperref[III.1.4.1]{1.4.1}). It is immediate that for every $f\in A_n$,
multiplication by $f$ defines a homomorphism
\[
C^\bullet(\mathfrak U,\mathcal F\otimes\mathcal L^{\otimes m})
\longrightarrow
C^\bullet(\mathfrak U,\mathcal F\otimes\mathcal L^{\otimes(m+n)}),
\]
whence a homomorphism
\[
H^\bullet(\mathfrak U,\mathcal F\otimes\mathcal L^{\otimes m})
\longrightarrow
H^\bullet(\mathfrak U,\mathcal F\otimes\mathcal L^{\otimes(m+n)}),
\]
which establishes the first assertion. On the other hand, for given
$f\in A_n$, it follows from (I, 9.3.2) that we have an isomorphism of
complexes of $(A_\bullet)_{(f)}$-modules
\[
C^\bullet(\mathfrak U\vert X_f,\mathcal F)
\xrightarrow{\sim}
\left(C^\bullet\left(\mathfrak U,
\bigoplus_{n\in\mathbb Z}\mathcal F\otimes\mathcal L^{\otimes n}
\right)\right)_{(f)}
\]
taking into account (I, 1.3.9, (ii)). Passing to the cohomology of these two
complexes, we deduce the isomorphism
(\hyperref[III.1.4.5.1]{1.4.5.1}), remembering that the functor
$M\mapsto M_{(f)}$ is exact on the category of graded
$A_\bullet$-modules.
\end{proof}

\begin{corollary}[1.4.6]
\label{III.1.4.6}
Suppose that the hypotheses of
(\hyperref[III.1.4.5]{1.4.5}) are satisfied, and suppose moreover that
$\mathcal L=\mathcal O_X$. If we set $A=\Gamma(X,\mathcal O_X)$, then,
for every $f\in A$, we have a canonical isomorphism
\[
H^\bullet(X_f,\mathcal F)
\xrightarrow{\sim}(H^\bullet(X,\mathcal F))_f
\]
of $A_f$-modules.
\end{corollary}

\begin{corollary}[1.4.7]
\label{III.1.4.7}
Let $X$ be a quasi-compact scheme and $f$ an element of
$\Gamma(X,\mathcal O_X)$.
\begin{enumerate}
\item[(i)] Suppose that the open subset $X_f$ is affine. Then, for every
quasi-coherent $\mathcal O_X$-module $\mathcal F$, every $i>0$ and every
$\xi\in H^i(X,\mathcal F)$, there exists an integer $n>0$ such that $f^n\xi=0$.
\item[(ii)] Conversely, suppose that $X_f$ is quasi-compact and that for
every quasi-coherent sheaf $\mathcal J$ of ideals of $\mathcal O_X$ and every
$\zeta\in H^1(X,\mathcal J)$, there exists $n>0$ such that $f^n\zeta=0$.
Then $X_f$ is affine.
\end{enumerate}
\end{corollary}

\begin{proof}
\begin{enumerate}
\item[(i)] If $X_f$ is affine, we have $H^i(X_f,\mathcal F)=0$ for every
$i>0$ (\hyperref[III.1.3.1]{1.3.1}), so the assertion follows directly
from (\hyperref[III.1.4.6]{1.4.6}).
\item[(ii)] By virtue of Serre's criterion (II, 5.2.1), it suffices to prove
that for every quasi-coherent ideal $\mathcal K$ of
$\mathcal O_X\vert X_f$, we have $H^1(X_f,\mathcal K)=0$. Since $X_f$ is a
quasi-compact open subset in a quasi-compact scheme $X$, there exists a
quasi-coherent ideal $\mathcal J$ of $\mathcal O_X$ such that
$\mathcal K=\mathcal J\vert X_f$ (I, 9.4.2). By virtue of
(\hyperref[III.1.4.6]{1.4.6}), we have
$H^1(X_f,\mathcal K)=(H^1(X,\mathcal J))_f$, and the hypothesis implies that
the right-hand side is zero, whence the conclusion.
\end{enumerate}
\end{proof}

\begin{remark}[1.4.8]
\label{III.1.4.8}
Note that (\hyperref[III.1.4.7]{1.4.7}, (i)) gives a simpler
proof of the relation (II, 4.5.13.2).
\end{remark}

\oldpage[III]{91}
\begin{lemma}[1.4.9]
\label{III.1.4.9}
Let $X$ be a quasi-compact scheme,
$\mathfrak U=(U_i)_{1\leq i\leq n}$ a finite covering of $X$ by affine
open subsets, and $\mathcal F$ a quasi-coherent $\mathcal O_X$-module. The
complex of sheaves $\mathcal C^\bullet(\mathfrak U,\mathcal F)$ defined
by the covering $\mathfrak U$ (G, II, 5.2) is then a quasi-coherent
$\mathcal O_X$-module.
\end{lemma}

\begin{proof}
It follows from the definitions (G, II, 5.2) that
$\mathcal C^p(\mathfrak U,\mathcal F)$ is the direct sum of the direct image
sheaves of the $\mathcal F\vert U_{i_0\ldots i_p}$ under the canonical
injection $U_{i_0\ldots i_p}\to X$. Now, the hypothesis that $X$ is a scheme
implies that these injections are affine morphisms (I, 5.5.6), so the
$\mathcal C^p(\mathfrak U,\mathcal F)$ are quasi-coherent (II, 1.2.6).
\end{proof}

\begin{proposition}[1.4.10]
\label{III.1.4.10}
Let $u:X\to Y$ be a separated and quasi-compact morphism. For every
quasi-coherent $\mathcal O_X$-module $\mathcal F$, the
$R^qu_*(\mathcal F)$ are quasi-coherent $\mathcal O_Y$-modules.
\end{proposition}

\begin{proof}
The question being local on $Y$, we may suppose $Y$ affine. Then $X$ is a
finite union of affine open subsets $U_i$ ($1\leq i\leq n$); let
$\mathfrak U$ be the covering $(U_i)$. Moreover, since $Y$ is a scheme,
it follows from (I, 5.5.10) that for every affine open subset $V\subset Y$,
the canonical injection $u^{-1}(V)\to X$ is an affine morphism; we
conclude from this ((\hyperref[III.1.4.1]{1.4.1}) and (G, II, 5.2)) that
we have a canonical isomorphism
\begin{equation}
\label{III.1.4.10.1}
\tag{1.4.10.1}
H^\bullet(u^{-1}(V),\mathcal F)
\xrightarrow{\sim}H^\bullet(\Gamma(V,\mathcal K^\bullet))
\end{equation}
where we have set
$\mathcal K^\bullet=u_*(\mathcal C^\bullet(\mathfrak U,\mathcal F))$. By
virtue of (\hyperref[III.1.4.9]{1.4.9}) and (I, 9.2.2),
$\mathcal K^\bullet$ is a quasi-coherent $\mathcal O_Y$-module; furthermore,
it constitutes a \emph{complex of sheaves}, since
$\mathcal C^\bullet(\mathfrak U,\mathcal F)$ does. It then follows from the
definition of the cohomology
$\mathcal H^\bullet(\mathcal K^\bullet)$ (G, II, 4.1) that the latter
consists of quasi-coherent $\mathcal O_Y$-modules (I, 4.1.1). Since (for $V$
affine in $Y$) the functor $\Gamma(V,\mathcal G)$ is exact in
$\mathcal G$ on the category of quasi-coherent $\mathcal O_Y$-modules,
we have (G, II, 4.1)
\begin{equation}
\label{III.1.4.10.2}
\tag{1.4.10.2}
H^\bullet(\Gamma(V,\mathcal K^\bullet))
=\Gamma(V,\mathcal H^\bullet(\mathcal K^\bullet)).
\end{equation}
Finally, note that it follows from the definition of the canonical
homomorphism
\[
H^\bullet(\mathfrak U,\mathcal F)\longrightarrow H^\bullet(X,\mathcal F)
\]
given in (G, II, 5.2) that if $V'\subset V$ is a second affine open subset
of $Y$, the diagram
\[
\xymatrix{
H^\bullet(u^{-1}(V),\mathcal F)\ar[r]^{\sim}\ar[d]
&H^\bullet(\Gamma(V,\mathcal K^\bullet))\ar[d]\\
H^\bullet(u^{-1}(V'),\mathcal F)\ar[r]^{\sim}
&H^\bullet(\Gamma(V',\mathcal K^\bullet))
}
\]
is commutative. We therefore conclude from the above that the isomorphisms
(\hyperref[III.1.4.10.1]{1.4.10.1}) define an isomorphism of
$\mathcal O_Y$-modules
\begin{equation}
\label{III.1.4.10.3}
\tag{1.4.10.3}
R^\bullet u_*(\mathcal F)
\xrightarrow{\sim}\mathcal H^\bullet(\mathcal K^\bullet)
\end{equation}
\oldpage[III]{92}
and consequently $R^\bullet u_*(\mathcal F)$ is quasi-coherent.
\end{proof}

It follows moreover from (\hyperref[III.1.4.10.3]{1.4.10.3}),
(\hyperref[III.1.4.10.2]{1.4.10.2}) and
(\hyperref[III.1.4.10.1]{1.4.10.1}) that:

\begin{corollary}[1.4.11]
\label{III.1.4.11}
Under the hypotheses of (\hyperref[III.1.4.10]{1.4.10}), for every
affine open subset $V$ of $Y$, the canonical homomorphism
\begin{equation}
\label{III.1.4.11.1}
\tag{1.4.11.1}
H^q(u^{-1}(V),\mathcal F)\longrightarrow
\Gamma(V,R^qu_*(\mathcal F))
\end{equation}
is an isomorphism for every $q\geq0$.
\end{corollary}

\begin{corollary}[1.4.12]
\label{III.1.4.12}
Suppose that the hypotheses of
(\hyperref[III.1.4.10]{1.4.10}) are satisfied, and suppose moreover that
$Y$ is quasi-compact. Then there exists an integer $r>0$ such that for every
quasi-coherent $\mathcal O_X$-module $\mathcal F$ and every integer $q>r$,
we have $R^qu_*(\mathcal F)=0$. If $Y$ is affine, we can take for $r$ an
integer such that there exists a covering of $X$ consisting of $r$ affine
open subsets.
\end{corollary}

\begin{proof}
Since $Y$ can be covered by finitely many affine open subsets, we are
reduced to proving the second assertion, by virtue of
(\hyperref[III.1.4.11]{1.4.11}). Now, if $\mathfrak U$ is a
covering of $X$ by $r$ affine open subsets, we have
$H^q(\mathfrak U,\mathcal F)=0$ for $q>r$, since the cochains of
$C^q(\mathfrak U,\mathcal F)$ are alternating; the conclusion therefore
follows from (\hyperref[III.1.4.1]{1.4.1}).
\end{proof}

\begin{corollary}[1.4.13]
\label{III.1.4.13}
Suppose that the hypotheses of
(\hyperref[III.1.4.10]{1.4.10}) are satisfied, and suppose moreover that
$Y=\operatorname{Spec}(A)$ is affine. Then, for every quasi-coherent
$\mathcal O_X$-module $\mathcal F$ and every $f\in A$, we have
\[
\Gamma(Y_f,R^qu_*(\mathcal F))
=(\Gamma(Y,R^qu_*(\mathcal F)))_f
\]
up to canonical isomorphism.
\end{corollary}

\begin{proof}
Indeed, this follows from the fact that $R^qu_*(\mathcal F)$ is a
quasi-coherent $\mathcal O_Y$-module (I, 1.3.7).
\end{proof}

\begin{proposition}[1.4.14]
\label{III.1.4.14}
Let $f:X\to Y$ be a separated quasi-compact morphism and $g:Y\to Z$ an
affine morphism. For every quasi-coherent $\mathcal O_X$-module
$\mathcal F$, the canonical homomorphism
$R^p(g\circ f)_*(\mathcal F)\to g_*(R^pf_*(\mathcal F))$ (0, 12.2.5.2) is
bijective for every $p$.
\end{proposition}

\begin{proof}
Indeed, for every affine open subset $W$ of $Z$, $g^{-1}(W)$ is an affine
open subset of $Y$. The homomorphism of presheaves defining the canonical
homomorphism
\[
R^p(g\circ f)_*(\mathcal F)\longrightarrow g_*(R^pf_*(\mathcal F))
\]
(0, 12.2.5) is therefore bijective by virtue of
(\hyperref[III.1.4.11]{1.4.11}).
\end{proof}

\begin{proposition}[1.4.15]
\label{III.1.4.15}
Let $u:X\to Y$ be a separated quasi-compact morphism and $v:Y'\to Y$ a
flat morphism of preschemes (0\textsubscript{I}, 6.7.1); let
$u'=u_{(Y')}$, so that we have the commutative diagram
\begin{equation}
\label{III.1.4.15.1}
\tag{1.4.15.1}
\xymatrix{
X\ar[d]_u & X'=X_{(Y')}\ar[l]_{v'}\ar[d]^{u'}\\
Y & Y'\ar[l]_v
}
\end{equation}
Then, for every quasi-coherent $\mathcal O_X$-module $\mathcal F$,
$R^qu_*'(\mathcal F')$, where
$\mathcal F'=v'^*(\mathcal F)=\mathcal F\otimes_{\mathcal O_Y}\mathcal O_{Y'}$,
is canonically isomorphic to
$R^qu_*(\mathcal F)\otimes_{\mathcal O_Y}\mathcal O_{Y'}
=v^*(R^qu_*(\mathcal F))$ for every $q\geq0$.
\end{proposition}

\oldpage[III]{93}
\begin{proof}
The canonical homomorphism
$\rho:\mathcal F\to v_*'(v'^*(\mathcal F))$ (0\textsubscript{I}, 4.4.3.2)
defines by functoriality a homomorphism
\begin{equation}
\label{III.1.4.15.2}
\tag{1.4.15.2}
R^qu_*(\mathcal F)\longrightarrow R^qu_*(v_*'(\mathcal F')).
\end{equation}
On the other hand, setting $w=u\circ v'=v\circ u'$, we have the
canonical homomorphisms (0, 12.2.5.1 and 12.2.5.2)
\begin{equation}
\label{III.1.4.15.3}
\tag{1.4.15.3}
R^qu_*(v_*'(\mathcal F'))\longrightarrow R^qw_*(\mathcal F')
\longrightarrow v_*(R^qu_*'(\mathcal F')).
\end{equation}
Composing (\hyperref[III.1.4.15.3]{1.4.15.3}) and
(\hyperref[III.1.4.15.2]{1.4.15.2}), we have a homomorphism
\[
\psi:R^qu_*(\mathcal F)\longrightarrow v_*(R^qu_*'(\mathcal F'))
\]
and finally we deduce from it the canonical homomorphism (whose definition
involves \emph{no hypothesis on $v$})
\begin{equation}
\label{III.1.4.15.4}
\tag{1.4.15.4}
\psi^\sharp:v^*(R^qu_*(\mathcal F))\longrightarrow R^qu_*'(\mathcal F')
\end{equation}
which we have to prove is an isomorphism when $v$ is
\emph{flat}. It is clear that the question is local on $Y$ and $Y'$, and we
can therefore suppose that $Y=\operatorname{Spec}(A)$,
$Y'=\operatorname{Spec}(B)$; we will moreover use the

\begin{lemma}[1.4.15.5]
\label{III.1.4.15.5}
Let $\varphi:A\to B$ be a ring homomorphism,
$Y=\operatorname{Spec}(A)$, $X=\operatorname{Spec}(B)$,
$f:X\to Y$ the morphism corresponding to $\varphi$, and $M$ a $B$-module.
For the $\mathcal O_X$-module $\widetilde M$ to be $f$-flat
(0\textsubscript{I}, 6.7.1), it is necessary and sufficient that $M$ be a
flat $A$-module. In particular, for the morphism $f$ to be flat, it is
necessary and sufficient that $B$ be a flat $A$-module.
\end{lemma}

Indeed, this follows from the definition (0\textsubscript{I}, 6.7.1) and
from (0\textsubscript{I}, 6.3.3), taking into account (I, 1.3.4).

This being so, it follows from
(\hyperref[III.1.4.11.1]{1.4.11.1}) and from the definitions of the
homomorphisms (\hyperref[III.1.4.15.3]{1.4.15.3})
(cf. 0, 12.2.5) that $\psi$ then corresponds to the composite morphism
\[
H^q(X,\mathcal F)\xrightarrow{\rho_q}
H^q(X,v_*'(v'^*(\mathcal F)))\xrightarrow{\theta_q}
H^q(X',v'^*(v_*'(v'^*(\mathcal F))))\xrightarrow{\sigma_q}
H^q(X',v'^*(\mathcal F))
\]
where $\rho_q$ and $\sigma_q$ are the homomorphisms corresponding in
cohomology to the canonical morphisms $\rho$ and
$\sigma:v'^*(v_*'(\mathcal G'))\to\mathcal G'$, and $\theta_q$ is the
$\varphi$-morphism (0, 12.1.3.1) relative to the $\mathcal O_X$-module
$v_*'(v'^*(\mathcal F))$. But by functoriality of $\theta_q$, we have the
commutative diagram
\[
\xymatrix{
H^q(X,\mathcal F)\ar[rr]^{\rho_q}\ar[dd]_{\theta_q}
&&H^q(X,v_*'(v'^*(\mathcal F)))\ar[dd]^{\theta_q}\\\\
H^q(X',v'^*(\mathcal F))\ar[rr]^{v'^*(\rho_q)}
&&H^q(X',v'^*(v_*'(v'^*(\mathcal F))))
}
\]
\oldpage[III]{94}
and since by definition (0\textsubscript{I}, 4.4.3) $v'^*(\rho)$ is
the inverse of $\sigma$, we see that the composite morphism considered above
is finally none other than $\theta_q$; $\psi^\sharp$ is consequently the
associated $B$-homomorphism
\[
H^q(X,\mathcal F)\otimes_A B\longrightarrow H^q(X',\mathcal F').
\]
Since $u$ is quasi-compact, $X$ is a finite union of affine open subsets $U_i$
($1\leq i\leq r$); let $\mathfrak U$ be the covering $(U_i)$. Furthermore,
$v$ is an affine morphism, hence so is $v'$ (II, 1.6.2,
(iii)), and the $U_i'=v'^{-1}(U_i)$ consequently form an affine open
covering $\mathfrak U'$ of $X'$. We then know (0, 12.1.4.2) that the diagram
\[
\xymatrix{
H^q(\mathfrak U,\mathcal F)\ar[r]^{\theta_q}\ar[d]
&H^q(\mathfrak U',\mathcal F')\ar[d]\\
H^q(X,\mathcal F)\ar[r]^{\theta_q}&H^q(X',\mathcal F')
}
\]
is commutative, and the vertical arrows are isomorphisms since
$X$ and $X'$ are schemes (\hyperref[III.1.4.1]{1.4.1}). It therefore
suffices to prove that the canonical $\varphi$-morphism
$\theta_q:H^q(\mathfrak U,\mathcal F)\to H^q(\mathfrak U',\mathcal F')$
is such that the associated $B$-homomorphism
\[
H^q(\mathfrak U,\mathcal F)\otimes_A B
\longrightarrow H^q(\mathfrak U',\mathcal F')
\]
is an isomorphism. For every sequence
$\mathbf s=(i_k)_{0\leq k\leq p}$ of $p+1$ indices of $[1,r]$, set
\[
U_{\mathbf s}=\bigcap_{k=0}^pU_{i_k},\qquad
U_{\mathbf s}'=\bigcap_{k=0}^pU_{i_k}'=v'^{-1}(U_{\mathbf s}),\qquad
M_{\mathbf s}=\Gamma(U_{\mathbf s},\mathcal F),\qquad
M_{\mathbf s}'=\Gamma(U_{\mathbf s}',\mathcal F').
\]
The canonical map
$M_{\mathbf s}\otimes_A B\to M_{\mathbf s}'$ is an isomorphism
(I, 1.6.5), so the canonical map
$C^p(\mathfrak U,\mathcal F)\otimes_A B\to
C^p(\mathfrak U',\mathcal F')$ is an isomorphism, by which
$d\otimes1$ is identified with the coboundary operator
$C^p(\mathfrak U',\mathcal F')\to C^{p+1}(\mathfrak U',\mathcal F')$.
Since $B$ is a \emph{flat} $A$-module, it follows immediately from the
definition of the cohomology modules that the canonical map
$H^q(\mathfrak U,\mathcal F)\otimes_A B\to
H^q(\mathfrak U',\mathcal F')$ is an isomorphism
(0\textsubscript{I}, 6.1.1). This result will be generalized later (\S~6).
\end{proof}

\begin{corollary}[1.4.16]
\label{III.1.4.16}
Let $A$ be a ring, $X$ an $A$-scheme of finite type, and $B$ an $A$-algebra
faithfully flat over $A$. For $X$ to be affine, it is necessary and sufficient
that $X\otimes_A B$ be affine.
\end{corollary}

\begin{proof}
The condition is evidently necessary (I, 3.2.2); let us show that it is
sufficient. Since $X$ is separated over $A$ and the morphism
$\operatorname{Spec}(B)\to\operatorname{Spec}(A)$ is flat, it follows from
(1.4.15) that we have
\begin{equation}
\label{III.1.4.16.1}
\tag{1.4.16.1}
H^i(X\otimes_A B,\mathcal F\otimes_A B)
=H^i(X,\mathcal F)\otimes_A B
\end{equation}
for every $i\geq0$ and every quasi-coherent $\mathcal O_X$-module $\mathcal F$.
If $X\otimes_A B$ is affine, the left-hand side of
(\hyperref[III.1.4.16.1]{1.4.16.1}) is zero for $i=1$, so the same holds
for $H^1(X,\mathcal F)$, since $B$ is a faithfully flat $A$-module.
Since $X$ is a quasi-compact scheme, we conclude by Serre's criterion
(II, 5.2.1).
\end{proof}

\begin{proposition}[1.4.17]
\label{III.1.4.17}
Let $X$ be a prescheme and
$0\to\mathcal F\xrightarrow{u}\mathcal G\xrightarrow{v}\mathcal H\to0$
an exact sequence of $\mathcal O_X$-modules. If $\mathcal F$ and $\mathcal H$
are quasi-coherent, so is $\mathcal G$.
\end{proposition}

\oldpage[III]{95}
\begin{proof}
The question being local on $X$, we may suppose
$X=\operatorname{Spec}(A)$ affine, and it then suffices to prove that
$\mathcal G$ satisfies conditions d~1) and d~2) of (I, 1.4.1) (with
$V=X$). The verification of d~2) is immediate, for if
$t\in\Gamma(X,\mathcal G)$ has zero restriction to $D(f)$, the same holds
for its image $v(t)\in\Gamma(X,\mathcal H)$; there is therefore $m>0$ such
that $f^mv(t)=v(f^mt)=0$ (I, 1.4.1), and since $\Gamma$ is left exact,
$f^mt=u(s)$, where $s\in\Gamma(X,\mathcal F)$; since $u$ is injective, the
restriction of $s$ to $D(f)$ is zero, whence (I, 1.4.1) the existence of an
integer $n>0$ such that $f^ns=0$; we finally deduce
$f^{m+n}t=u(f^ns)=0$.

Let us now verify d~1), and let $t'\in\Gamma(D(f),\mathcal G)$; since
$\mathcal H$ is quasi-coherent, there exists an integer $m$ such that
$f^mv(t')=v(f^mt')$ extends to a section
$z\in\Gamma(X,\mathcal H)$ (I, 1.4.1). But by virtue of
(\hyperref[III.1.3.1]{1.3.1}) (or of (I, 5.1.9.2)) applied to the
quasi-coherent $\mathcal O_X$-module $\mathcal F$, the sequence
$\Gamma(X,\mathcal G)\to\Gamma(X,\mathcal H)\to0$ is exact, so there exists
$t\in\Gamma(X,\mathcal G)$ such that $z=v(t)$; we therefore see that
$v(f^mt'-t'')=0$, denoting by $t''$ the restriction of $t$ to $D(f)$;
hence $f^mt'-t''=u(s')$, where $s'\in\Gamma(D(f),\mathcal F)$. But since
$\mathcal F$ is quasi-coherent, there exists an integer $n>0$ such that $f^ns'$
extends to a section $s\in\Gamma(X,\mathcal F)$; since
$f^{m+n}t'-f^nt''=u(f^ns')$, we see that $f^{m+n}t'$ is the restriction to
$D(f)$ of the section $f^nt+u(s)\in\Gamma(X,\mathcal G)$, which completes
the proof.
\end{proof}

